\section{Objetivo del Informe}

El objetivo central del presente informe es evaluar de manera cuantitativa y cualitativa los resultados obtenidos tras la implementación del sistema ERP Odoo y el despliegue del ecosistema omnicanal en Textiles y Confecciones Atlas E.I.R.L. La finalidad de esta evaluación es validar cómo la transformación digital logró erradicar la dependencia de procesos manuales e informales, dotando a la empresa de una infraestructura tecnológica segura, centralizada y escalable.

Asimismo, el informe busca documentar y corroborar el cumplimiento de las métricas de éxito (KPIs) establecidas al inicio del proyecto. Específicamente, la evaluación se centrará en evidenciar los siguientes logros:

\begin{itemize}
    \item \textbf{Eficiencia en la Atención a Provincias:} La reducción drástica del tiempo promedio de respuesta a consultas de stock para clientes a distancia, pasando de un rango inicial de 45 a 60 minutos (debido a la verificación física) a \textbf{menos de 5 minutos}, gracias a la consulta inmediata en la base de datos de Odoo.
    \item \textbf{Exactitud del Inventario:} La evolución desde un 0\% de exactitud digital hacia un \textbf{75\% de coincidencia inicial} en el cierre de stock, estableciendo las bases para alcanzar a mediano plazo un 95\% de precisión en tiempo real.
    \item \textbf{Reducción de Errores Operativos:} La consecución de una reducción del \textbf{15\% en los errores de inventario} y discrepancias de facturación en el mostrador durante el primer semestre de uso del sistema integrado.
    \item \textbf{Expansión Omnicanal:} El salto de poseer 0 canales oficiales a consolidar \textbf{2 canales digitales activos} con identidad corporativa profesional, logrando transicionar de un esquema reactivo a la recepción sostenida de más de 5 consultas estructuradas por mes.
\end{itemize}

A través del análisis detallado de estos indicadores, el documento confirmará el impacto directo de la digitalización en la optimización de los recursos, la mejora en el servicio al cliente y la recuperación del rol estratégico de la gerencia.
